% The rule given to the model for the traditional-pedagogy reading.
% Reproduced verbatim from labelling/targets/passes/
% trad_pedagogy_RULE_traditional_pedagogy.md -- do not reword.

\begin{quote}
\footnotesize

\textbf{The question.} You are reading a crawl of a school's own public
website: typically the homepage, an about or welcome page, and a values or
ethos page. Three pages is all there is. \emph{Does this school's public
self-presentation signal a traditional pedagogical identity?} Traditional
pedagogy here means teaching understood as the transmission of knowledge by a
subject expert to a class: the teacher instructs from the front, the
curriculum is a body of knowledge to be learned, and the school says so as
part of what it is. Answer YES or NO. There is no middle band.

\textbf{What this is not.} It is not a judgement about the school; traditional
is neither better nor worse than the alternative. It is not about how the
school actually teaches, only what it says about itself. It is not about
strictness, behaviour, uniform or discipline. It is not about academic
ambition: ``high expectations'', ``excellence'' and ``aspiration'' appear on
most school websites and say nothing about pedagogy. It is not about a school
being old or having a long history.

\textbf{The test that does the work.} A naming is not an account. The question
is whether the school has said something about how teaching works, not whether
it has used a word from a vocabulary list. If the adjectives are stripped out,
has this school said anything about what happens in its classrooms?

YES looks like a school that describes teaching as instruction by subject
experts; that says knowledge comes first and skills follow from it; that
describes whole-class teaching, direct instruction, practice and retrieval, or
a knowledge-rich curriculum \emph{and} says what it means by that; that names a
pedagogical authority or doctrine and stands by it; or that calls its own
teaching traditional and elaborates.

NO looks like a school that uses one of those words in passing and never
returns to it; that says it has ``subject specialists'' or ``expert teachers'',
which nearly every secondary school does; that says it is ``knowledge-rich''
once, in a list of other things; or that uses ``traditional'' about something
that is not teaching.

\textbf{The four traps.} Each is a reason to answer NO.
(1)~``Subject specialists'' or ``expert teachers'': a description of who is
employed, not of how they teach.
(2)~A doctrine word used once, in passing: the word is the start of the
evidence, not the evidence.
(3)~``Traditional'' applied to something that is not pedagogy: traditional
values, uniform, expectations of behaviour, or a traditional Key~Stage~4
curriculum, which is a fact about curriculum structure.
(4)~Generic high-standards language: rigour, challenge, stretch, aiming high.
This is ambition, not pedagogy.

\textbf{The other direction.} Some schools signal the opposite: child-led
enquiry, personalised learning journeys, learning styles, the teacher as
facilitator, skills before knowledge, or project-based learning as the
organising idea. If a progressive framework is the school's primary account of
teaching, the answer is NO, even if a traditional-sounding word appears
somewhere. A school with no pedagogical claims at all is also NO.

\textbf{How to answer.} Borderline goes NO. If it could be argued either way on
what is in front of you, the answer is NO; most schools are NO. Record a
verdict, a confidence, and for a YES the exact sentence from the document that
carries it. If an exact sentence cannot be copied, the answer is NO. Judge each
document only on what is in it, using nothing known about the school from
elsewhere, and do not compare documents with each other.

\end{quote}
